\chapter{系统架构与工程实现}
\label{chap:implementation}

\cref{chap:design}从方法层说明了本文如何通过 Agent 自主任务调度方法、原子检测能力工具化和 LLM 分层信息整合策略解决建筑幕墙韧性评估中的任务选择与信息融合问题。本章转向工程链路实现，重点说明上述方法如何落地为一个可运行、可部署、可观测、可恢复的软件系统。系统实现的核心不是简单把若干模型串联起来，而是围绕项目、图像、任务、产物、复核和报告构建完整的工程链路，使 Agent 能力在受控状态机中稳定发挥作用。

\section{关键术语}
\label{sec:implementation-key-terms}

本章涉及的对象更多偏向工程实现，包括服务模块、通信方式、存储对象和运维机制等。为便于阅读，现对本章关键术语进行说明。

\begin{table}[!htb]
  \caption{系统架构与工程实现关键术语}\label{tab:implementation-key-terms}
  \centering
  \begin{tabular}{@{}p{0.16\linewidth}p{0.72\linewidth}@{}} \toprule
    \textbf{术语} & \textbf{含义} \\ \midrule
    PSM & 服务进程或模块的逻辑名称，用于区分 Core API、Core BIZ、Activity 服务、CronJob 和日志采集等运行单元。 \\
    Core API & 面向前端提供 HTTP、鉴权入口和 WebSocket 通道等的接入服务。 \\
    Core BIZ & 系统业务中心，负责状态机推进、数据库事务、对象存储授权、Worker 编排、回调处理、事件发布和定时任务等。 \\
    Activity 服务 & 独立智能能力服务，包括 Classification、Reasoning 和 Summary，分别封装分类 Agent、原子检测模型和总结 Agent。 \\
    Worker & 后台任务执行器，异步推进检测任务和报告任务，避免用户请求长时间阻塞。 \\
    回调 & Activity 服务完成任务后，通过 gRPC 将状态、结果、产物清单和错误信息上报给 Core BIZ 的机制。 \\
    RocketMQ 事件 & Core BIZ 在图像、检测任务或报告状态变化后发布的项目事件，由 Core API 消费并转发给前端。 \\
    WebSocket Scope & WebSocket 推送范围，包括图像资产、检测任务和报告三类，用于隔离不同阶段的实时消息。 \\
    预签名 URL & Core BIZ 为前端或 Activity 服务生成的有时效对象存储访问地址，用于上传或下载图像与报告源文件。 \\
    产物 & Reasoning 服务生成并上传到 MinIO 的中间图像，如标注图、区域图、浮层图、梯度图、直线图和频谱图等。 \\
    CronJob & Core BIZ 内部定时任务，用于上传超时处理和 MinIO 脏对象清理等后台维护工作。 \\ \bottomrule
  \end{tabular}
\end{table}

\section{总体架构设计与实现}
\label{sec:implementation-overall-architecture-design}

系统采用前后端分离与服务化架构。前端 Core Web 负责项目工作台、项目阶段页面、图像资产管理、检测进度展示、检测结果查看、人工复核和报告展示；后端按照职责拆分为 Core API、Core BIZ、Activity Classification、Activity Reasoning、Activity Summary 和 ICW Common；基础设施包含 MySQL、Redis、MinIO 和 RocketMQ；部署层由 Docker Compose 统一组织容器、网络、环境变量、日志采集和依赖启动顺序。

\begin{figure}[!htb]
  \centering
  \resizebox{\textwidth}{!}{%
    \begin{tikzpicture}[
      font=\scriptsize,
      region/.style={draw, rounded corners=4pt, line width=0.65pt, align=center},
      layerTitle/.style={font=\bfseries\scriptsize, anchor=west, align=left},
      cell/.style={draw, rounded corners=2pt, align=center, minimum height=0.46cm, line width=0.42pt, inner sep=2pt},
      userbox/.style={region, draw=blue!55!black, fill=blue!4},
      corebox/.style={region, draw=teal!55!black, fill=teal!5},
      activitybox/.style={region, draw=orange!75!black, fill=orange!7},
      infrabox/.style={region, draw=gray!70!black, fill=gray!7},
      deploybox/.style={region, draw=purple!55!black, fill=purple!5},
      usercell/.style={cell, draw=blue!55!black, fill=blue!8},
      corecell/.style={cell, draw=teal!55!black, fill=teal!10},
      activitycell/.style={cell, draw=orange!75!black, fill=orange!12},
      infracell/.style={cell, draw=gray!70!black, fill=gray!12},
      deploycell/.style={cell, draw=purple!55!black, fill=purple!9}
    ]
      \node[deploybox, minimum width=14.30cm, minimum height=1.30cm] at (0, 5.45) {};
      \node[layerTitle] at (-7.00, 5.77) {部署与运行治理层};
      \node[deploycell, minimum width=1.80cm] at (-5.15, 5.24) {Docker};
      \node[deploycell, minimum width=1.80cm] at (-3.09, 5.24) {统一网络};
      \node[deploycell, minimum width=1.80cm] at (-1.03, 5.24) {环境变量};
      \node[deploycell, minimum width=1.80cm] at (1.03, 5.24) {依赖启动};
      \node[deploycell, minimum width=1.80cm] at (3.09, 5.24) {日志采集};
      \node[deploycell, minimum width=1.80cm] at (5.15, 5.24) {镜像发布};

      \node[userbox, minimum width=14.30cm, minimum height=1.55cm] at (0, 3.83) {};
      \node[layerTitle] at (-7.00, 4.28) {用户与前端层};
      \node[usercell, minimum width=1.80cm] at (-5.15, 3.64) {浏览器用户\\项目操作};
      \node[usercell, minimum width=1.80cm] at (-3.09, 3.64) {Core Web\\React 页面};
      \node[usercell, minimum width=1.80cm] at (-1.03, 3.64) {阶段流程\\状态展示};
      \node[usercell, minimum width=1.80cm] at (1.03, 3.64) {图像结果\\可视化};
      \node[usercell, minimum width=1.80cm] at (3.09, 3.64) {人工复核\\报告打印};
      \node[usercell, minimum width=1.80cm] at (5.15, 3.64) {WebSocket\\实时刷新};

      \node[corebox, minimum width=14.30cm, minimum height=2.35cm] at (0, 1.72) {};
      \node[layerTitle] at (-7.00, 2.56) {核心业务服务层};
      \node[corecell, minimum width=2.10cm] at (-5.00, 1.96) {Core API\\HTTP 接入};
      \node[corecell, minimum width=2.10cm] at (-2.50, 1.96) {鉴权权限\\请求日志};
      \node[corecell, minimum width=2.10cm] at (0.00, 1.96) {WebSocket\\Hub};
      \node[corecell, minimum width=2.10cm] at (2.50, 1.96) {Core BIZ\\状态机中心};
      \node[corecell, minimum width=2.10cm] at (5.00, 1.96) {Worker\\任务编排};
      \node[corecell, minimum width=2.10cm] at (-5.00, 1.04) {ICW Common\\Proto / Enum};
      \node[corecell, minimum width=2.10cm] at (-2.50, 1.04) {Repository\\事务封装};
      \node[corecell, minimum width=2.10cm] at (0.00, 1.04) {回调处理\\聚合推进};
      \node[corecell, minimum width=2.10cm] at (2.50, 1.04) {事件生产\\状态广播};
      \node[corecell, minimum width=2.10cm] at (5.00, 1.04) {CronJob\\后台维护};

      \node[activitybox, minimum width=14.30cm, minimum height=2.35cm] at (0, -0.81) {};
      \node[layerTitle] at (-7.00, 0.03) {智能活动服务层};
      \node[activitycell, minimum width=2.50cm] at (-4.80, -0.57) {Classification\\分类 Agent};
      \node[activitycell, minimum width=2.50cm] at (-1.60, -0.57) {Reasoning\\原子检测能力};
      \node[activitycell, minimum width=2.50cm] at (1.60, -0.57) {Summary\\总结 Agent};
      \node[activitycell, minimum width=2.50cm] at (4.80, -0.57) {Python Worker\\模型脚本};
      \node[activitycell, minimum width=2.50cm] at (-4.80, -1.49) {图像附件\\任务代码};
      \node[activitycell, minimum width=2.50cm] at (-1.60, -1.49) {检测产物与指标\\report.json};
      \node[activitycell, minimum width=2.50cm] at (1.60, -1.49) {报告上下文\\source.txt};
      \node[activitycell, minimum width=2.50cm] at (4.80, -1.49) {Summary 输出\\压缩 JSON};

      \node[infrabox, minimum width=14.30cm, minimum height=1.55cm] at (0, -2.94) {};
      \node[layerTitle] at (-7.00, -2.50) {基础设施与数据层};
      \node[infracell, minimum width=1.80cm] at (-5.15, -3.14) {MySQL\\结构化状态};
      \node[infracell, minimum width=1.80cm] at (-3.09, -3.14) {Redis\\令牌与缓存};
      \node[infracell, minimum width=1.80cm] at (-1.03, -3.14) {MinIO\\图像与产物};
      \node[infracell, minimum width=1.80cm] at (1.03, -3.14) {RocketMQ\\项目事件};
      \node[infracell, minimum width=1.80cm] at (3.09, -3.14) {Log Sidecar\\分 PSM 日志};
      \node[infracell, minimum width=1.80cm] at (5.15, -3.14) {Config\\环境变量配置};
    \end{tikzpicture}
  }
  \caption{系统总体架构与服务关系}\label{fig:implementation-overall-architecture}
\end{figure}

如\cref{fig:implementation-overall-architecture} 所示，系统总体架构可以概括为五个层次：部署与运行治理层覆盖所有服务，提供容器编排、网络隔离、环境配置、日志采集和镜像发布能力；用户与前端层负责页面交互和实时状态呈现；核心业务服务层负责接入、鉴权、状态机、任务编排、回调处理和事件生产；智能活动服务层负责分类、检测和总结等相对独立的智能能力；基础设施与数据层负责结构化状态、缓存、对象产物、事件消息和日志留存。各服务或模块的实现职责见\cref{tab:service-division}。

\begin{table}[!htb]
  \caption{系统服务划分与实现职责}\label{tab:service-division}
  \centering
  \begin{tabular}{@{}p{0.18\linewidth}p{0.76\linewidth}@{}} \toprule
    \textbf{服务或模块} & \textbf{实现职责} \\ \midrule
    Core Web & React 前端页面、阶段化项目流程、图像和检测结果可视化、人工复核、报告展示与打印 \\
    Core API & HTTP 路由、鉴权中间件、请求日志、跨域处理、WebSocket 连接、RocketMQ 事件消费与前端广播 \\
    Core BIZ & 业务状态机、数据库事务、对象存储封装、Redis 缓存、检测 Worker、报告 Worker、回调处理、事件生产和 CronJob \\
    Activity Classification & 调用分类 Agent，上传图像附件，解析受控 JSON 输出，回调检测任务代码 \\
    Activity Reasoning & 统一封装五类 Python 原子检测能力，负责原图下载、模型执行、产物上传和结果回调 \\
    Activity Summary & 调用图像总结 Agent 和项目报告 Agent，压缩 Agent 输出 JSON，并回调总结与报告结果 \\
    ICW Common & Protocol Buffers 协议、共享枚举、gRPC 调用封装、错误码、日志工具、JSON 归一化和通用工具函数 \\ \bottomrule
  \end{tabular}
\end{table}

该架构遵循“智能能力外置，业务状态内聚”的原则。Classification、Reasoning 和 Summary 可以独立演进各自 Agent 或模型逻辑，但不直接拥有业务状态；项目、图像、任务、复核和报告状态集中由 Core BIZ 控制；Core API 不承担复杂业务事务，只负责接入与事件转发；Core Web 只展示和触发业务动作，不自行判断最终状态。这样的边界划分避免了 Agent 与模型服务越权，也使失败恢复和日志追踪更清晰。

从工程复杂度看，用户一次“启动 Agent 智能检测”背后会触发多条异步链路：Core Web 发起 HTTP 请求，Core API 鉴权并转发到 Core BIZ，Core BIZ 为图像创建检测主任务并投递 Worker，Worker 调用分类 Agent，分类回调后创建多个子任务，Reasoning 服务执行 Python 模型并上传产物，Summary 服务生成单图总结，BIZ 发布 RocketMQ 事件，API 消费后通过 WebSocket 推送给前端。系统必须保证这些链路在并发和失败情况下仍能保持状态一致。

\section{通信与协议设计}
\label{sec:implementation-communication-and-protocol-design}

系统通信分为四类。第一类是前端到 Core API 的 HTTP 请求，用于登录、项目管理、图像上传授权、检测启动、结果查询、复核更新和报告查询。第二类是 Core API、Core BIZ 与 Activity 服务之间的 gRPC 请求，用于服务间强类型调用。第三类是 Core BIZ 到 Activity 服务的异步启动与 Activity 服务到 Core BIZ 的回调，用于长任务结果上报。第四类是 Core BIZ 发布 RocketMQ 项目事件、Core API 消费事件并通过 WebSocket 广播给前端，用于实时状态刷新。

Protocol Buffers 是内部协议的核心。系统将 Auth、ProjectCore、ProjectProfile、ProjectAssets、ProjectDetection、ProjectReview、ProjectReport、Socket 和 User 等 BIZ 接口定义在 Common 仓库中；Activity 侧定义 Classification、Reasoning 和 Summary 三类服务接口；前端 API 类型也由 proto 生成 TypeScript 文件。通过共享协议，项目进度、图像状态、检测任务代码、主任务状态、子任务状态、复核结论、报告状态和 WebSocket Scope 等均使用统一枚举，以避免服务间常量漂移。

请求链路中还通过请求 ID 实现日志关联。Core API 在 HTTP 中间件生成请求 ID，gRPC 调用时通过 \texttt{Metadata} 透传，Activity 服务异步执行和回调时继续携带该请求 ID。这样可以在日志中关联一次用户操作、BIZ 处理、Agent 调用、模型执行和回调结果，对长链路排障具有直接价值。

RocketMQ 用于承接业务状态变化后的异步项目事件。Core BIZ 在图像上传、检测主任务、检测子任务、单图总结和项目报告等状态发生变化后，将项目 ID、事件类型和事件载荷发布到统一 Topic；Core API 作为消费者订阅该 Topic，解析事件类型并映射为前端需要的 WebSocket Scope。通过消息队列，业务状态推进与前端连接管理被解耦：BIZ 只负责在事务完成后发布事实事件，不直接感知在线客户端；API 只负责消费事件和广播消息，不参与业务事务。该机制既避免长任务回调直接阻塞 HTTP 请求，也使图像资产、检测进度和报告状态能够在页面中以增量方式刷新。

WebSocket 设计采用“票据 + Scope”的方式。前端先通过鉴权 HTTP 请求创建连接票据，再使用项目 ID、Scope 和 Ticket 建立 WebSocket。Scope 将图像资产、检测任务和报告状态隔离，避免所有页面都接收全量事件。API 侧 Hub 按项目和 Scope 管理连接，RocketMQ 消费到项目事件后只广播给对应项目与对应 Scope 的客户端。

\section{模块设计与实现}
\label{sec:implementation-module-design-and-implementation}

本节说明各服务模块的核心设计与实现，每个模块都给出职责说明和核心时序图。

\subsection{Core API 服务}
\label{subsec:implementation-api-service}

Core API 是前端访问系统的统一入口。该服务基于 Gin 注册 HTTP 路由，覆盖认证、用户头像、项目核心流程、项目基础信息、图像资产、智能检测、人工复核、评估报告和 WebSocket 票据等接口。API 服务通过鉴权中间件从令牌中解析用户身份，通过项目访问中间件校验用户是否拥有项目权限，然后将请求转换为 Core BIZ gRPC 请求。

除 HTTP 接入外，Core API 还承担实时事件转发职责。Core BIZ 不直接维护前端连接，而是将图像状态、检测任务状态和报告状态写入 RocketMQ。Core API 作为消费者订阅项目事件，根据事件类型选择 WebSocket Scope，再向对应项目连接广播消息。这样设计使 BIZ 不依赖 WebSocket 实现，也使 API 服务成为前端通信的唯一入口。

如\cref{fig:implementation-api-service-sequence} 所示，Core API 的实现重点在于把同步查询和异步事件统一到前端体验中。首次进入页面或手动刷新时，前端通过 HTTP 获取全量状态；任务执行过程中，API 通过 WebSocket 推送增量事件。前端收到事件后静默刷新局部数据，避免页面反复进入加载态。

\begin{figure}[!htb]
  \centering
  \resizebox{\linewidth}{!}{%
    \begin{tikzpicture}[
      font=\scriptsize,
      >=stealth,
      participant/.style={draw, rounded corners=2pt, align=center, minimum width=1.9cm, minimum height=0.62cm, fill=blue!6, line width=0.55pt},
      service/.style={participant},
      infra/.style={participant},
      lifeline/.style={dashed, gray!70, line width=0.45pt},
      activation/.style={draw, fill=gray!18, minimum width=0.18cm, inner sep=0pt, line width=0.45pt},
      msg/.style={->, line width=0.55pt},
      ret/.style={->, dashed, line width=0.55pt},
      label/.style={font=\tiny, align=center, inner sep=1pt}
    ]
      \node[participant] (web) at (0, 0) {Core Web};
      \node[service] (api) at (2.25, 0) {Core API\\HTTP};
      \node[service] (auth) at (4.50, 0) {鉴权与\\项目权限};
      \node[service] (biz) at (6.75, 0) {Core BIZ\\gRPC};
      \node[infra] (mq) at (9.00, 0) {RocketMQ};
      \node[service] (hub) at (11.25, 0) {WebSocket\\Hub};

      \foreach \x in {0,2.25,4.50,6.75,9.00,11.25} {
        \draw[lifeline] (\x,-0.45) -- (\x,-9.15);
      }
      \node[activation, minimum height=5.75cm] at (2.25,-3.775) {};
      \node[activation, minimum height=0.60cm] at (4.50,-3.15) {};
      \node[activation, minimum height=0.60cm] at (6.75,-4.45) {};
      \node[activation, minimum height=0.70cm] at (9.00,-7.70) {};
      \node[activation, minimum height=0.70cm] at (11.25,-8.40) {};

      \draw[msg] (0,-0.90) -- node[label, above] {HTTP 请求 / 携带 token 与 project\_id} (2.17,-0.90);
      \draw[msg] (2.33,-1.55) -- (3.12,-1.55) -- node[label, right, pos=0.5] {路由匹配 / 绑定请求参数 / 生成 request id} (3.12,-2.15) -- (2.33,-2.15);
      \draw[msg] (2.33,-2.85) -- node[label, above] {解析用户身份 / 校验项目权限} (4.42,-2.85);
      \draw[ret] (4.42,-3.45) -- node[label, below] {user\_id / project 权限结果} (2.33,-3.45);
      \draw[msg] (2.33,-4.15) -- node[label, above] {封装 gRPC payload / metadata 透传 request id} (6.67,-4.15);
      \draw[ret] (6.67,-4.75) -- node[label, below] {业务响应 / 业务错误码} (2.33,-4.75);
      \draw[msg] (2.33,-5.45) -- (3.12,-5.45) -- node[label, right, pos=0.5] {统一 ApiEnvelope / 错误映射 / 访问日志} (3.12,-6.05) -- (2.33,-6.05);
      \draw[ret] (2.17,-6.65) -- node[label, below] {HTTP 响应} (0,-6.65);
      \draw[msg] (6.83,-7.35) -- node[label, above] {业务状态变化后发布项目事件} (8.92,-7.35);
      \draw[msg] (9.08,-8.05) -- node[label, above] {API 消费事件并解析 event\_type / scope} (11.17,-8.05);
      \draw[msg] (11.17,-8.75) -- node[label, above] {按 project\_id + scope 广播} (0,-8.75);
    \end{tikzpicture}
  }
  \caption{Core API 服务核心时序图}\label{fig:implementation-api-service-sequence}
\end{figure}

\subsection{Core BIZ 服务}
\label{subsec:implementation-biz-service}

Core BIZ 是系统中最复杂的服务，也是工程链路的状态中心。它负责项目创建、项目阶段推进、图像组管理、图像上传授权、图像上报、检测任务创建、人工复核、报告任务创建、对象存储访问、Redis 缓存、RocketMQ 事件生产、CronJob 和后台 Worker。所有会改变业务状态的操作都集中在 BIZ 服务中执行。

BIZ 服务的核心设计是“服务层 + Repository 层 + Worker 层”。服务层实现 gRPC 接口和业务规则，Repository 层封装 MySQL、Redis、MinIO 和 RocketMQ 操作，Worker 层异步推进检测和报告任务。Activity 服务回调 BIZ 后，BIZ 在数据库事务中写入结果、判断聚合状态、创建下一阶段任务并发布项目事件。

如\cref{fig:implementation-biz-service-sequence} 所示，BIZ 服务既是请求处理者，也是长任务编排者。Worker 只负责启动 Activity 服务，真正结果通过回调回到 BIZ。这样可以避免 Worker 阻塞等待模型推理，也能保证所有状态变更都回到同一业务中心处理。

\begin{figure}[!htb]
  \centering
  \resizebox{\linewidth}{!}{%
    \begin{tikzpicture}[
      font=\scriptsize,
      >=stealth,
      participant/.style={draw, rounded corners=2pt, align=center, minimum width=1.75cm, minimum height=0.62cm, fill=blue!6, line width=0.55pt},
      service/.style={participant},
      infra/.style={participant},
      lifeline/.style={dashed, gray!70, line width=0.45pt},
      activation/.style={draw, fill=gray!18, minimum width=0.18cm, inner sep=0pt, line width=0.45pt},
      msg/.style={->, line width=0.55pt},
      ret/.style={->, dashed, line width=0.55pt},
      label/.style={font=\tiny, align=center, inner sep=1pt}
    ]
      \node[participant] (api) at (0, 0) {Core API};
      \node[service] (biz) at (2.10, 0) {Core BIZ\\Service};
      \node[infra] (db) at (4.20, 0) {MySQL};
      \node[infra] (store) at (6.30, 0) {MinIO / Redis};
      \node[service] (worker) at (8.40, 0) {Worker};
      \node[service] (activity) at (10.50, 0) {Activity\\服务};
      \node[infra] (mq) at (12.60, 0) {RocketMQ};

      \foreach \x in {0,2.10,4.20,6.30,8.40,10.50,12.60} {
        \draw[lifeline] (\x,-0.45) -- (\x,-11.25);
      }
      \node[activation, minimum height=9.95cm] at (2.10,-5.875) {};
      \node[activation, minimum height=0.60cm] at (4.20,-3.15) {};
      \node[activation, minimum height=0.60cm] at (4.20,-7.85) {};
      \node[activation, minimum height=0.60cm] at (6.30,-4.45) {};
      \node[activation, minimum height=0.60cm] at (6.30,-9.15) {};
      \node[activation, minimum height=1.40cm] at (8.40,-6.15) {};
      \node[activation, minimum height=0.70cm] at (10.50,-6.50) {};
      \node[activation, minimum height=0.40cm] at (12.60,-10.35) {};

      \draw[msg] (0,-0.90) -- node[label, above] {启动检测 / 推进阶段 / 查询状态} (2.02,-0.90);
      \draw[msg] (2.18,-1.55) -- (2.95,-1.55) -- node[label, right, pos=0.5] {校验项目阶段 / 幂等检查 / 构造领域对象} (2.95,-2.15) -- (2.18,-2.15);
      \draw[msg] (2.18,-2.85) -- node[label, above] {开启事务：写入项目、图像、检测主任务} (4.12,-2.85);
      \draw[ret] (4.12,-3.45) -- node[label, below] {任务 UUID / 当前状态 / 版本字段} (2.18,-3.45);
      \draw[msg] (2.18,-4.15) -- node[label, above] {生成预签名 URL / 写入 Redis 临时控制} (6.22,-4.15);
      \draw[ret] (6.22,-4.75) -- node[label, below] {对象地址 / 缓存结果} (2.18,-4.75);
      \draw[msg] (2.18,-5.45) -- node[label, above] {提交后台任务并释放用户请求} (8.32,-5.45);
      \draw[msg] (8.48,-6.15) -- node[label, above] {按状态启动分类、检测、总结或报告} (10.42,-6.15);
      \draw[ret] (10.42,-6.85) -- node[label, below] {Activity 回调结果 / 失败原因 / 产物清单} (2.18,-6.85);
      \draw[msg] (2.18,-7.55) -- node[label, above] {回调事务：写结果、聚合子状态、创建下一阶段任务} (4.12,-7.55);
      \draw[ret] (4.12,-8.15) -- node[label, below] {落库成功 / 回滚失败} (2.18,-8.15);
      \draw[msg] (2.18,-8.85) -- node[label, above] {读取/写入 source.txt、产物授权、URL 缓存} (6.22,-8.85);
      \draw[ret] (6.22,-9.45) -- node[label, below] {对象 Key / 缓存命中 / 授权策略} (2.18,-9.45);
      \draw[msg] (2.18,-10.15) -- node[label, above] {发布图像、检测或报告状态事件} (12.52,-10.15);
      \draw[ret] (2.02,-10.85) -- node[label, below] {gRPC 响应或异步处理结束} (0,-10.85);
    \end{tikzpicture}
  }
  \caption{Core BIZ 服务核心时序图}\label{fig:implementation-biz-service-sequence}
\end{figure}

Core BIZ 内部还包含两个 CronJob。上传中图像超时任务定期扫描长时间未完成上报的图像，将其状态置为失败并发布事件；MinIO 脏对象清理任务每小时扫描对象存储，删除失败图像、失败子任务、失败报告和数据库已无引用对象对应的文件夹，并设置 30 分钟保护窗口，避免刚生成但尚未完成数据库状态同步的对象被误删。

\subsection{Classification 服务}
\label{subsec:implementation-classification-service}

Classification 服务封装分类 Agent。BIZ 启动分类任务时传入任务 UUID、图像 UUID 和原图预签名下载地址。Classification 服务下载并压缩图像后，通过 Agent 平台文件上传接口上传附件，再通过流式对话接口获取模型输出。服务端将 Agent 输出归一化为 JSON 对象，解析 \texttt{task\_codes}，转换为共享枚举后回调 BIZ。

\cref{fig:implementation-classification-service-sequence} 的关键点是：分类 Agent 只决定任务代码，不创建子任务。输出异常时，Classification 服务回调失败状态；输出成功时，BIZ 仍要进行枚举转换、任务去重和子任务创建。这样既利用 Agent 对幕墙材质的语义判断能力，又避免 Agent 越过业务状态机。

\begin{figure}[!htb]
  \centering
  \resizebox{\linewidth}{!}{%
    \begin{tikzpicture}[
      font=\scriptsize,
      >=stealth,
      participant/.style={draw, rounded corners=2pt, align=center, minimum width=1.95cm, minimum height=0.62cm, fill=blue!6, line width=0.55pt},
      infra/.style={participant},
      agent/.style={participant},
      lifeline/.style={dashed, gray!70, line width=0.45pt},
      activation/.style={draw, fill=gray!18, minimum width=0.18cm, inner sep=0pt, line width=0.45pt},
      msg/.style={->, line width=0.55pt},
      ret/.style={->, dashed, line width=0.55pt},
      label/.style={font=\tiny, align=center, inner sep=1pt}
    ]
      \node[participant] (biz) at (0, 0) {Core BIZ};
      \node[participant] (cls) at (2.55, 0) {Classification\\Service};
      \node[infra] (minio) at (5.10, 0) {MinIO\\原图};
      \node[agent] (file) at (7.65, 0) {Agent\\文件上传};
      \node[agent] (chat) at (10.20, 0) {分类 Agent\\对话};

      \foreach \x in {0,2.55,5.10,7.65,10.20} {
        \draw[lifeline] (\x,-0.45) -- (\x,-11.10);
      }
      \node[activation, minimum height=9.85cm] at (2.55,-5.825) {};
      \node[activation, minimum height=0.60cm] at (5.10,-3.15) {};
      \node[activation, minimum height=0.60cm] at (7.65,-5.75) {};
      \node[activation, minimum height=0.70cm] at (10.20,-7.10) {};

      \draw[msg] (0,-0.90) -- node[label, above] {StartClassification(task\_uuid, image\_uuid, url)} (2.47,-0.90);
      \draw[msg] (2.63,-1.55) -- (3.48,-1.55) -- node[label, right, pos=0.5] {加载 prompt / 校验 bot id 与 token / 创建请求上下文} (3.48,-2.15) -- (2.63,-2.15);
      \draw[msg] (2.63,-2.85) -- node[label, above] {按预签名 URL 下载原图附件} (5.02,-2.85);
      \draw[ret] (5.02,-3.45) -- node[label, below] {图像字节流} (2.63,-3.45);
      \draw[msg] (2.63,-4.15) -- (3.48,-4.15) -- node[label, right, pos=0.5] {压缩图像 / 控制附件体积 / 生成文件名} (3.48,-4.75) -- (2.63,-4.75);
      \draw[msg] (2.63,-5.45) -- node[label, above] {上传图像附件} (7.57,-5.45);
      \draw[ret] (7.57,-6.05) -- node[label, below] {file\_id} (2.63,-6.05);
      \draw[msg] (2.63,-6.75) -- node[label, above] {提示词 + file\_id 流式调用分类 Agent} (10.12,-6.75);
      \draw[ret] (10.12,-7.45) -- node[label, below] {模型输出：JSON 字符串} (2.63,-7.45);
      \draw[msg] (2.63,-8.15) -- (3.48,-8.15) -- node[label, right, pos=0.5] {CompactAgentJSONObjectString / 解析 task\_codes} (3.48,-8.75) -- (2.63,-8.75);
      \draw[msg] (2.63,-9.45) -- (3.48,-9.45) -- node[label, right, pos=0.5] {非法代码拦截 / 去重 / 固定顺序排序 / 空集合允许} (3.48,-10.05) -- (2.63,-10.05);
      \draw[ret] (2.47,-10.75) -- node[label, below] {ReportClassificationResult(task\_codes / error)} (0,-10.75);
    \end{tikzpicture}
  }
  \caption{Classification 服务核心时序图}\label{fig:implementation-classification-service-sequence}
\end{figure}

\subsection{Reasoning 服务}
\label{subsec:implementation-reasoning-service}

Reasoning 服务是 Go 工程层与 Python 模型层之间的桥梁。服务启动时注册五类原子检测能力，并根据任务代码检查模型目录中对应权重文件是否存在、是否可读取、是否为 Git LFS 指针文件。任务执行时，Reasoning 服务从 BIZ 获取最新原图下载地址，下载到隔离运行目录，再调用常驻 Python Worker 执行对应检测模块。模型执行结束后，Go 层读取压缩后的 \texttt{report.json}，再从 BIZ 获取产物上传授权，将目录中的产物并发上传到 MinIO，并带着结果 JSON 和产物清单回调 BIZ。

\begin{figure}[!htb]
  \centering
  \resizebox{\linewidth}{!}{%
    \begin{tikzpicture}[
      font=\scriptsize,
      >=stealth,
      participant/.style={draw, rounded corners=2pt, align=center, minimum width=1.85cm, minimum height=0.62cm, fill=blue!6, line width=0.55pt},
      infra/.style={participant},
      model/.style={participant},
      lifeline/.style={dashed, gray!70, line width=0.45pt},
      activation/.style={draw, fill=gray!18, minimum width=0.18cm, inner sep=0pt, line width=0.45pt},
      msg/.style={->, line width=0.55pt},
      ret/.style={->, dashed, line width=0.55pt},
      label/.style={font=\tiny, align=center, inner sep=1pt}
    ]
      \node[participant] (biz) at (0, 0) {Core BIZ};
      \node[participant] (rsn) at (2.15, 0) {Reasoning\\Service};
      \node[participant] (registry) at (4.30, 0) {能力注册表};
      \node[infra] (minio) at (6.45, 0) {MinIO};
      \node[model] (py) at (8.60, 0) {Python\\Worker};
      \node[model] (model) at (10.75, 0) {检测模型\\脚本};

      \foreach \x in {0,2.15,4.30,6.45,8.60,10.75} {
        \draw[lifeline] (\x,-0.45) -- (\x,-13.75);
      }
      \node[activation, minimum height=12.45cm] at (2.15,-7.125) {};
      \node[activation, minimum height=0.60cm] at (4.30,-1.90) {};
      \node[activation, minimum height=0.60cm] at (0,-3.20) {};
      \node[activation, minimum height=0.60cm] at (6.45,-5.80) {};
      \node[activation, minimum height=2.00cm] at (8.60,-7.80) {};
      \node[activation, minimum height=0.60cm] at (10.75,-7.80) {};
      \node[activation, minimum height=0.60cm] at (6.45,-11.10) {};

      \draw[msg] (0,-0.90) -- node[label, above] {StartReasoning(task\_code, image\_uuid)} (2.07,-0.90);
      \draw[msg] (2.23,-1.60) -- node[label, above] {按 task\_code 查询执行器与产物规则} (4.22,-1.60);
      \draw[ret] (4.22,-2.20) -- node[label, below] {模型入口 / 动态产物清单} (2.23,-2.20);
      \draw[msg] (2.07,-2.90) -- node[label, above] {获取最新原图 URL 与上传授权} (0.08,-2.90);
      \draw[ret] (0.08,-3.50) -- node[label, below] {presign get / artifact policy} (2.07,-3.50);
      \draw[msg] (2.23,-4.20) -- (3.05,-4.20) -- node[label, right, pos=0.5] {创建隔离运行目录 / 初始化输入与产物路径} (3.05,-4.80) -- (2.23,-4.80);
      \draw[msg] (2.23,-5.50) -- node[label, above] {下载原图到运行目录} (6.37,-5.50);
      \draw[ret] (6.37,-6.10) -- node[label, below] {原图文件} (2.23,-6.10);
      \draw[msg] (2.23,-6.80) -- node[label, above] {提交 Python 检测任务} (8.52,-6.80);
      \draw[msg] (8.68,-7.50) -- node[label, above] {懒加载模型并执行脚本} (10.67,-7.50);
      \draw[ret] (10.67,-8.10) -- node[label, below] {report.json / 图像产物} (8.68,-8.10);
      \draw[ret] (8.52,-8.80) -- node[label, below] {检测结果目录} (2.23,-8.80);
      \draw[msg] (2.23,-9.50) -- (3.05,-9.50) -- node[label, right, pos=0.5] {读取并压缩 report.json / 计算产物 SHA256} (3.05,-10.10) -- (2.23,-10.10);
      \draw[msg] (2.23,-10.80) -- node[label, above] {上传全部产物并校验 Uploaded=true} (6.37,-10.80);
      \draw[ret] (6.37,-11.40) -- node[label, below] {object key / sha256 / uploaded} (2.23,-11.40);
      \draw[msg] (2.23,-12.10) -- (3.05,-12.10) -- node[label, right, pos=0.5] {产物数量与动态清单校验 / 失败则阻断成功回调} (3.05,-12.70) -- (2.23,-12.70);
      \draw[ret] (2.07,-13.35) -- node[label, below] {ReportReasoningResult(result JSON / artifacts / error)} (0,-13.35);
    \end{tikzpicture}
  }
  \caption{Reasoning 服务核心时序图}\label{fig:implementation-reasoning-service-sequence}
\end{figure}

如\cref{fig:implementation-reasoning-service-sequence} 所示，Reasoning 服务不依赖 BIZ 在启动时传入可能过期的原图下载地址和产物上传策略，而是在真正执行时向 BIZ 获取最新授权。这解决了长队列场景中预签名 URL 过期导致下载或上传失败的问题。对于除玻璃爆裂以外的检测任务，Reasoning 服务要求产物不为空且全部上传成功，否则回调失败，不将不完整检测结果视为成功。

Python 模型采用懒加载方式。Go 服务通过常驻 Worker 向 Python 发送任务，Python 首次使用某类检测模型时加载权重，后续复用已加载模型。这样既保留 Python 图像处理和深度学习生态，又避免每次检测重复加载模型造成高延迟。

\subsection{Summary 服务}
\label{subsec:implementation-summary-service}

\begin{figure}[!htb]
  \centering
  \resizebox{\linewidth}{!}{%
    \begin{tikzpicture}[
      font=\scriptsize,
      >=stealth,
      participant/.style={draw, rounded corners=2pt, align=center, minimum width=1.95cm, minimum height=0.62cm, fill=blue!6, line width=0.55pt},
      infra/.style={participant},
      agent/.style={participant},
      lifeline/.style={dashed, gray!70, line width=0.45pt},
      activation/.style={draw, fill=gray!18, minimum width=0.18cm, inner sep=0pt, line width=0.45pt},
      msg/.style={->, line width=0.55pt},
      ret/.style={->, dashed, line width=0.55pt},
      label/.style={font=\tiny, align=center, inner sep=1pt}
    ]
      \node[participant] (biz) at (0, 0) {Core BIZ};
      \node[participant] (summary) at (2.40, 0) {Summary\\Service};
      \node[infra] (minio) at (4.80, 0) {MinIO\\source.txt};
      \node[agent] (file) at (7.20, 0) {Agent\\附件上传};
      \node[agent] (agent) at (9.60, 0) {总结 Agent\\对话};

      \foreach \x in {0,2.40,4.80,7.20,9.60} {
        \draw[lifeline] (\x,-0.45) -- (\x,-11.20);
      }
      \node[activation, minimum height=9.90cm] at (2.40,-5.85) {};
      \node[activation, minimum height=0.60cm] at (4.80,-4.50) {};
      \node[activation, minimum height=0.60cm] at (7.20,-5.80) {};
      \node[activation, minimum height=0.70cm] at (9.60,-7.15) {};

      \draw[msg] (0,-0.90) -- node[label, above] {StartDetectionSummary / StartProjectSummary} (2.32,-0.90);
      \draw[msg] (2.48,-1.60) -- (3.35,-1.60) -- node[label, right, pos=0.5] {读取任务类型 / 加载 detection 或 project prompt} (3.35,-2.20) -- (2.48,-2.20);
      \draw[msg] (2.48,-2.90) -- (3.35,-2.90) -- node[label, right, pos=0.5] {检测总结：使用结构化 JSON 文本输入} (3.35,-3.50) -- (2.48,-3.50);
      \draw[msg] (2.48,-4.20) -- node[label, above] {项目报告：下载 source.txt} (4.72,-4.20);
      \draw[ret] (4.72,-4.80) -- node[label, below] {结构化源数据文本} (2.48,-4.80);
      \draw[msg] (2.48,-5.50) -- node[label, above] {上传 source.txt 附件或复用单图文本} (7.12,-5.50);
      \draw[ret] (7.12,-6.10) -- node[label, below] {file\_id / 输入上下文} (2.48,-6.10);
      \draw[msg] (2.48,-6.80) -- node[label, above] {提示词 + 文本 / 附件调用 Agent} (9.52,-6.80);
      \draw[ret] (9.52,-7.50) -- node[label, below] {检测总结 JSON 或项目报告 JSON} (2.48,-7.50);
      \draw[msg] (2.48,-8.20) -- (3.35,-8.20) -- node[label, right, pos=0.5] {CompactAgentJSONObjectString / JSON object 校验} (3.35,-8.80) -- (2.48,-8.80);
      \draw[msg] (2.48,-9.50) -- (3.35,-9.50) -- node[label, right, pos=0.5] {项目报告校验 group\_count、image\_count 与源数据一致} (3.35,-10.10) -- (2.48,-10.10);
      \draw[ret] (2.32,-10.80) -- node[label, below] {ReportDetectionSummaryResult / ReportProjectSummaryResult} (0,-10.80);
    \end{tikzpicture}
  }
  \caption{Summary 服务核心时序图}\label{fig:implementation-summary-service-sequence}
\end{figure}

Summary 服务包含图像总结和项目报告两类任务。如\cref{fig:implementation-summary-service-sequence} 所示，图像总结任务接收 BIZ 组装的单图检测结果 JSON，直接调用检测总结 Agent，并要求输出固定七字段 JSON：总体概述、五类缺陷总结和后续建议。项目报告任务接收 BIZ 提供的 \texttt{source.txt} 预签名地址，先下载该附件，再将附件上传至报告 Agent，要求输出固定项目报告 JSON。

Summary 服务的实现重点在于输出稳定性。Agent 输出需要经过 JSON 压缩和对象校验，项目报告还要校验输出中的图像组数量和图像数量是否与源数据一致，防止 Agent 生成与项目数据不匹配的报告。数据库中保存的是压缩 JSON 字符串，前端按字段解析和展示。

\subsection{Core Web 服务}
\label{subsec:implementation-web-service}

\begin{figure}[!htb]
  \centering
  \resizebox{\linewidth}{!}{%
    \begin{tikzpicture}[
      font=\scriptsize,
      >=stealth,
      participant/.style={draw, rounded corners=2pt, align=center, minimum width=1.85cm, minimum height=0.62cm, fill=blue!6, line width=0.55pt},
      service/.style={participant},
      infra/.style={participant},
      lifeline/.style={dashed, gray!70, line width=0.45pt},
      activation/.style={draw, fill=gray!18, minimum width=0.18cm, inner sep=0pt, line width=0.45pt},
      msg/.style={->, line width=0.55pt},
      ret/.style={->, dashed, line width=0.55pt},
      label/.style={font=\tiny, align=center, inner sep=1pt}
    ]
      \node[participant] (user) at (0, 0) {用户};
      \node[service] (page) at (2.10, 0) {阶段页面\\组件};
      \node[service] (store) at (4.20, 0) {状态与\\弹窗组件};
      \node[service] (http) at (6.30, 0) {HTTP\\Client};
      \node[service] (api) at (8.40, 0) {Core API};
      \node[infra] (ws) at (10.50, 0) {WebSocket};

      \foreach \x in {0,2.10,4.20,6.30,8.40,10.50} {
        \draw[lifeline] (\x,-0.45) -- (\x,-11.50);
      }
      \node[activation, minimum height=8.90cm] at (2.10,-5.35) {};
      \node[activation, minimum height=9.50cm] at (4.20,-6.35) {};
      \node[activation, minimum height=6.10cm] at (6.30,-5.35) {};
      \node[activation, minimum height=0.60cm] at (8.40,-4.60) {};
      \node[activation, minimum height=0.70cm] at (10.50,-6.65) {};

      \draw[msg] (0,-0.90) -- node[label, above] {进入项目阶段页面} (2.02,-0.90);
      \draw[msg] (2.18,-1.60) -- node[label, above] {初始化页面状态 / 显示首次加载态} (4.12,-1.60);
      \draw[msg] (2.18,-2.30) -- node[label, above] {查询基础信息、资产、检测、报告} (6.22,-2.30);
      \draw[msg] (6.38,-3.00) -- (7.20,-3.00) -- node[label, right, pos=0.5] {请求前判断 access token 是否需要刷新} (7.20,-3.60) -- (6.38,-3.60);
      \draw[msg] (6.38,-4.30) -- node[label, above] {HTTP 请求 / payload request 格式} (8.32,-4.30);
      \draw[ret] (8.32,-4.90) -- node[label, below] {结构化响应 / 401 则排队刷新 token} (6.38,-4.90);
      \draw[ret] (6.22,-5.60) -- node[label, below] {页面数据} (2.18,-5.60);
      \draw[msg] (2.18,-6.30) -- node[label, above] {根据阶段 scope 创建 socket ticket 并建连} (10.42,-6.30);
      \draw[ret] (10.42,-7.00) -- node[label, below] {图像 / 检测 / 报告状态事件} (2.18,-7.00);
      \draw[msg] (2.18,-7.70) -- node[label, above] {静默刷新受影响列表，不触发全屏加载} (6.22,-7.70);
      \draw[ret] (6.22,-8.40) -- node[label, below] {最新列表数据} (4.28,-8.40);
      \draw[msg] (0,-9.10) -- node[label, above] {查看进度 / 结果 / 复核 / 打印报告} (2.02,-9.10);
      \draw[msg] (2.18,-9.80) -- node[label, above] {打开弹窗，默认 Agent 总结页并重置 tab} (4.12,-9.80);
      \draw[msg] (4.28,-10.50) -- (5.12,-10.50) -- node[label, right, pos=0.5] {按生成类型和枚举字段渲染指标、产物、评论} (5.12,-11.10) -- (4.28,-11.10);
    \end{tikzpicture}
  }
  \caption{Core Web 服务核心时序图}\label{fig:implementation-web-service-sequence}
\end{figure}

Core Web 使用 React 和 TypeScript 实现。前端页面按项目阶段组织，包括项目基础信息、图像资产构建、Agent 智能检测、人工复核确认和评估报告生成。前端类型尽量复用 Common 中生成的 TypeScript 结构和枚举，避免手写常量与后端不一致。HTTP 请求统一经过 Axios 封装，处理令牌刷新、错误提示和 API 地址配置。

如\cref{fig:implementation-web-service-sequence} 所示，前端采用“首次全量加载 + WebSocket 增量事件 + 静默刷新”的方式维护页面状态。首次进入页面时显示加载态并拉取完整数据；WebSocket 连接建立或收到事件后，不再强制展示全页面加载，而是静默拉取最新列表并更新局部状态。检测结果弹窗、进度弹窗和报告页面均基于结构化数据渲染，而不是直接展示原始 JSON。

\subsection{ICW Common 模块}
\label{subsec:implementation-common-module}

ICW Common 是跨仓库共享基础设施。它包含 Protocol Buffers 定义、生成后的 Golang/TypeScript 代码、共享枚举、错误码、日志常量、JSON 压缩归一化、gRPC 调用封装、表格化日志输出和 ID 编解码等工具。Common 模块的作用是确保多个服务在协议和基础行为上保持一致。

\begin{figure}[!htb]
  \centering
  \resizebox{\linewidth}{!}{%
    \begin{tikzpicture}[
      font=\scriptsize,
      >=stealth,
      participant/.style={draw, rounded corners=2pt, align=center, minimum width=1.9cm, minimum height=0.62cm, fill=blue!6, line width=0.55pt},
      service/.style={participant},
      lifeline/.style={dashed, gray!70, line width=0.45pt},
      activation/.style={draw, fill=gray!18, minimum width=0.18cm, inner sep=0pt, line width=0.45pt},
      msg/.style={->, line width=0.55pt},
      ret/.style={->, dashed, line width=0.55pt},
      label/.style={font=\tiny, align=center, inner sep=1pt}
    ]
      \node[participant] (dev) at (0, 0) {开发者};
      \node[participant] (proto) at (2.30, 0) {Common\\Proto};
      \node[participant] (gen) at (4.60, 0) {buf / protoc\\生成器};
      \node[service] (go) at (6.90, 0) {Go 服务\\API/BIZ/Activity};
      \node[service] (web) at (9.20, 0) {Core Web\\TypeScript};
      \node[service] (runtime) at (11.50, 0) {运行期\\gRPC/HTTP};

      \foreach \x in {0,2.30,4.60,6.90,9.20,11.50} {
        \draw[lifeline] (\x,-0.45) -- (\x,-10.00);
      }
      \node[activation, minimum height=4.10cm] at (2.30,-2.95) {};
      \node[activation, minimum height=0.70cm] at (4.60,-1.95) {};
      \node[activation, minimum height=2.00cm] at (6.90,-5.30) {};
      \node[activation, minimum height=0.70cm] at (6.90,-9.25) {};
      \node[activation, minimum height=4.60cm] at (9.20,-7.30) {};
      \node[activation, minimum height=0.60cm] at (11.50,-8.60) {};

      \draw[msg] (0,-0.90) -- node[label, above] {维护接口、枚举、消息结构} (2.22,-0.90);
      \draw[msg] (2.38,-1.60) -- node[label, above] {执行 buf generate / protoc} (4.52,-1.60);
      \draw[ret] (4.52,-2.30) -- node[label, below] {Go pb、gRPC stub、TypeScript gen} (2.38,-2.30);
      \draw[msg] (2.38,-3.00) -- (3.20,-3.00) -- node[label, right, pos=0.5] {统一枚举：任务代码、状态、Scope、事件类型} (3.20,-3.60) -- (2.38,-3.60);
      \draw[msg] (2.38,-4.30) -- node[label, above] {发布 gRPC 接口、错误码、JSON 工具、日志工具} (6.82,-4.30);
      \draw[msg] (2.38,-5.00) -- node[label, above] {同步请求响应结构、枚举与常量} (9.12,-5.00);
      \draw[msg] (6.98,-5.70) -- (7.80,-5.70) -- node[label, right, pos=0.5] {go fmt / goimports / go vet / make verify} (7.80,-6.30) -- (6.98,-6.30);
      \draw[msg] (9.28,-7.00) -- (10.10,-7.00) -- node[label, right, pos=0.5] {src/gen 禁止 lint，业务代码直接引用生成类型} (10.10,-7.60) -- (9.28,-7.60);
      \draw[msg] (9.28,-8.30) -- node[label, above] {按生成类型封装 payload} (11.42,-8.30);
      \draw[msg] (11.42,-8.90) -- node[label, above] {HTTP / gRPC 调用进入后端} (6.98,-8.90);
      \draw[ret] (6.98,-9.60) -- node[label, below] {统一错误码、状态枚举和 JSON 归一化结果} (9.28,-9.60);
    \end{tikzpicture}
  }
  \caption{ICW Common 模块核心时序图}\label{fig:implementation-common-module-sequence}
\end{figure}

\cref{fig:implementation-common-module-sequence} 展示了 Common 模块的开发与使用链路。共享枚举对本文系统尤其重要，因为 Agent 输出的任务代码、WebSocket Scope、项目事件类型和检测状态都需要在前后端与多个后端服务间保持一致。统一协议减少了重复定义，也降低了后续新增检测能力时的维护成本。

\section{数据实体设计与实现}
\label{sec:implementation-data-entity-design-and-implementation}

系统数据层需要同时支撑项目阶段流转、图像资产管理、Agent 任务调度、原子检测结果、人工复核、项目报告和对象产物追踪。为避免将所有业务信息堆叠在单个 JSON 字段中，本文采用“结构化数据库 + Redis 临时状态 + MinIO 对象存储”的组合方式：MySQL 负责保存可查询、可聚合、可事务化的数据实体；Redis 负责保存令牌、缓存和短时控制状态；MinIO 负责保存图像、检测产物和报告源数据等二进制或大体积文件。三者分工明确，使系统既能支持复杂业务查询，又能控制数据库体积和对象访问权限。

\subsection{数据库设计与实现}
\label{subsec:implementation-database-design-and-implementation}

数据库设计围绕“项目级资产”和“图像级任务”展开。项目表保存项目名称、建筑名称、位置、建成年份、建筑描述、已知问题、评估目标、项目状态和项目进度；图像组表保存项目内图像分组名称、用户归属和排序信息，用于表达立面、区域或用户自定义图像集合；图像表保存单张图像的 UUID、所属图像组、文件名、尺寸、大小、上传状态和元数据，是后续检测主任务、子检测任务和人工复核的图像级数据入口；检测主任务表保存每张图像的检测整体状态、五类子任务是否执行、子任务 ID、图像总结任务 ID 和人工复核字段；五张子任务表分别保存不同原子检测能力的专属指标；图像总结任务表保存单图总结；项目报告表保存项目级报告 JSON 和状态。

本文系统的核心数据库实体及其作用见\cref{tab:implementation-core-entities}。

数据库事务用于保证状态推进的一致性。例如，分类回调成功后，BIZ 需要在同一事务内写入任务路由字段、创建子任务、关联子任务 ID、创建必要的总结任务并推进主任务状态；子任务回调后，BIZ 需要写入专属指标、保存产物 SHA256、更新子任务状态，并判断同一主任务下所有必要子任务是否完成；项目报告回调后，BIZ 需要写入报告 JSON、更新报告状态并发布报告事件。事务边界保证了前端查询到的状态与任务记录保持一致。

\begin{table}[!htb]
  \caption{核心数据库实体与作用}\label{tab:implementation-core-entities}
  \centering
  \begin{tabular}{@{}p{0.32\linewidth}p{0.62\linewidth}@{}} \toprule
    \textbf{数据实体} & \textbf{作用} \\ \midrule
    \texttt{projects} & 保存项目基础信息、项目状态和阶段进度，是所有业务数据的归属边界 \\
    \texttt{project\_groups} & 保存项目下的图像组名称、用户归属、项目归属和排序优先级，用于组织立面、区域或用户自定义图像集合 \\
    \texttt{project\_group\_images} & 保存单张图像的 UUID、所属图像组、文件名、MIME 类型、尺寸、大小、元数据、上传状态和上传时间 \\
    \texttt{project\_detection\_tasks} & 保存单张图像检测主任务、任务路由结果、主任务状态、总结任务和人工复核字段 \\
    五类检测子任务表 & 分别保存锈蚀、裂缝、污渍、平整度、爆裂检测结果、区域 JSON、产物 SHA256 和运行耗时 \\
    \texttt{project\_detection\_summary\_tasks} & 保存图像级总结任务状态和 Agent 输出结果 \\
    \texttt{project\_reports} & 保存项目级报告状态和项目报告 JSON \\ \bottomrule
  \end{tabular}
\end{table}

不同子任务表采用“公共任务字段 + 专属指标字段”的结构。公共字段便于统一查询和状态展示，专属字段便于前端按任务类型展示不同指标，也方便项目报告源数据按任务类型组装。相比所有检测结果都以一个大 JSON 字段保存，该结构更利于状态聚合、索引优化和工程维护。

\subsection{缓存设计与实现}
\label{subsec:implementation-cache-design-and-implementation}

缓存主要服务于三类短时数据。第一类是用户安全相关状态，包括邮箱验证码、登录失败计数、账号短时锁定、Access Token 黑名单和 Refresh Token 刷新互斥锁。邮箱验证码经过场景区分并在验证成功后消费，避免同一验证码重复使用；登录失败计数用于限制密码或验证码暴力尝试；刷新互斥锁用于防止同一 Refresh Token 被并发请求同时刷新。第二类是对象访问缓存，包括图像原图、缩略图、检测产物和项目缩略图的预签名下载 URL。由于前端查看检测结果和图像资产时会频繁访问同一对象，缓存预签名 URL 可以减少重复生成 URL 的开销。第三类是临时控制状态，例如 Socket Ticket 和上传中的图像状态保护。

Redis 中保存的内容均具有明确 TTL，不作为长期业务事实来源。业务事实仍以 MySQL 为准，Redis 只承担加速访问、短时互斥和临时安全控制。删除图像、重试检测或清理对象时，系统会同步清理相关预签名 URL 缓存，避免继续使用旧对象地址。

\subsection{对象存储设计与实现}
\label{subsec:implementation-object-storage-design-and-implementation}

系统使用 MinIO 保存项目缩略图、图像原图、图像缩略图、检测产物和项目报告源数据。数据库不保存二进制文件，只保存元数据、状态和必要的产物校验信息。前端上传图像时，BIZ 生成预签名 UPLOAD URL，前端将原图和缩略图直接上传到对象存储，再调用图像上报接口确认状态。检测产物上传时，Reasoning 服务向 BIZ 获取受限前缀上传授权，在按项目、图像和任务组织的路径下上传动态数量产物。

产物完整性是 Reasoning 链路的关键约束。Reasoning 服务上传产物后会返回名称、对象 Key、SHA256 和上传结果。BIZ 在回调处理阶段校验产物清单，对于需要产物的任务，必须保证所有产物 \texttt{uploaded=true} 且 SHA256 非空，才将子任务视为成功。否则该子任务失败，不写入成功结果，后续可以通过重试重新生成并上传产物。这样可以避免“检测结果成功但产物缺失”的页面异常。

对象存储与缓存层配合使用。由于图像和产物可能被前端反复查看，系统会将部分预签名下载 URL 缓存在 Redis 中，减少重复生成 URL 的开销。删除图像或清理脏对象时，BIZ 会同步清理对应对象缓存，避免前端继续使用已经失效的对象地址。

\section{系统状态机设计与实现}
\label{sec:implementation-system-state-machine-design-and-implementation}

状态机设计是系统保证长链路一致性的核心。如前文\cref{fig:requirements-stage-state-machine} 所示，系统中的状态机可以分为项目工作流主状态机、图像资产构建状态机、检测主任务状态机、检测子任务状态机和评估报告生成状态机。它们共同约束用户从项目创建、图像上传、Agent 智能检测、人工复核到评估报告生成的完整流程，避免前端直接跳转阶段或后端任务乱序推进。

项目阶段进度从项目初始化、基础信息完成、图像资产完成、检测完成、复核完成推进到项目完成。每一次阶段推进都由 Core BIZ 检查前置条件：基础信息阶段要求保存项目资料，图像资产阶段要求至少存在一张上传成功图像，检测阶段要求主任务全部完成，复核阶段要求用户完成确认，报告阶段只有报告成功生成后才允许将项目视为完成。前端按钮状态只作为交互约束，最终阶段判断始终由后端执行。

检测主任务状态机是任务编排的核心。主任务状态包括 \texttt{pending}、\texttt{classifying}、\texttt{detecting}、\texttt{summarizing}、\texttt{succeeded} 和 \texttt{failed}。分类 Agent 成功后，BIZ 根据任务代码创建必要子任务；子任务全部成功后创建单图总结；单图总结成功后主任务进入成功状态。任一必要环节失败时，主任务进入失败状态，并保留失败原因供前端展示和用户重试。

子任务状态包括 \texttt{pending}、\texttt{succeeded} 和 \texttt{failed}。Reasoning 回调成功不等于子任务必然成功，BIZ 还需要校验产物完整性、产物上传状态和任务指标合法性。报告生成状态同样包括 \texttt{pending}、\texttt{succeeded} 和 \texttt{failed}。报告生成失败时前端显示失败重试按钮，报告生成中允许刷新当前状态，报告生成成功后不再显示重试或刷新操作。通过这些状态机，系统把多服务异步执行转化为可查询、可重试和可解释的业务过程。

\section{安全性设计与实现}
\label{sec:implementation-security-design-and-implementation}

安全性设计覆盖用户身份管理、登录鉴权、项目权限、令牌生命周期、对象存储访问、WebSocket 连接、密码存储、服务边界和 Agent 输出约束等。系统的安全基本原则是：前端只保存必要的短期访问凭据，不暴露长期后端凭证；所有项目级数据访问均回到后端校验；对象存储访问通过短时授权完成；Agent 与模型服务只返回结构化结果，不直接写库或推进状态。

用户认证采用 Access Token 与 Refresh Token 组合。Access Token 是短时 JWT，包含用户 ID、邮箱、用户名、签发时间、过期时间、签发方和 jti，并使用 HMAC 签名算法；Refresh Token 是长期随机令牌，明文只返回给客户端，数据库只保存其哈希值和 Token ID。API 鉴权中间件只校验 Access Token，当前端检测到访问令牌即将过期或收到 401 后，再通过 \texttt{/auth/refresh} 使用 HttpOnly Cookie 中的 Refresh Token 换取新令牌。

\begin{figure}[!htb]
  \centering
  \resizebox{\linewidth}{!}{%
    \begin{tikzpicture}[
      font=\scriptsize,
      >=stealth,
      normal/.style={draw=blue!70!black, fill=blue!6, rounded corners=3pt, align=center, minimum width=3.05cm, minimum height=0.78cm, line width=0.55pt},
      success/.style={draw=green!55!black, fill=green!8, rounded corners=3pt, align=center, minimum width=3.05cm, minimum height=0.78cm, line width=0.55pt},
      fail/.style={draw=red!70!black, fill=red!7, rounded corners=3pt, align=center, minimum width=3.05cm, minimum height=0.78cm, line width=0.55pt},
      arrow/.style={->, line width=0.58pt},
      lane/.style={font=\bfseries\scriptsize, align=center, anchor=center},
      every node/.style={outer sep=0pt}
    ]
      \path[use as bounding box] (-9.25,-3.30) rectangle (9.25,3.05);
      \draw[draw=gray!70!black, fill=gray!7, rounded corners=4pt, line width=0.65pt]
        (-9.25,-3.30) rectangle (9.25,3.05);

      \node[lane] at (-8.29,2.25) {登录签发};
      \node[normal] (login) at (-5.95,2.25) {用户登录\\密码或邮箱验证码};
      \node[normal] (verify) at (-2.65,2.25) {BIZ 校验身份\\失败计数与锁定};
      \node[normal] (issue) at (0.65,2.25) {签发 Access Token\\生成 Refresh Token};
      \node[normal] (save) at (3.95,2.25) {MySQL 保存\\Refresh Token Hash};
      \node[normal] (cookie) at (7.25,2.25) {API 写入\\HttpOnly Cookie};

      \node[lane] at (-8.29,0.65) {访问校验};
      \node[normal] (request) at (-5.95,0.65) {前端请求 API\\携带 Access Token};
      \node[normal] (auth) at (-2.65,0.65) {API 调用 Me\\校验 JWT};
      \node[success] (valid) at (0.65,0.65) {Access Token\\有效};
      \node[success] (biz) at (3.95,0.65) {BIZ 返回用户\\继续业务请求};
      \node[fail] (expired) at (-2.65,-0.95) {Access Token\\过期或无效};

      \node[success] (refresh) at (-5.95,-0.95) {/auth/refresh\\读取 Cookie};
      \node[lane] at (-8.29,-2.55) {刷新轮换};
      \node[success] (lock) at (-5.95,-2.55) {Redis 设置刷新锁\\防止并发刷新};
      \node[success] (hash) at (-2.65,-2.55) {按 Token ID + Hash\\查询 Refresh Token};
      \node[normal] (rotate) at (0.65,-2.55) {轮换新令牌\\吊销旧 Refresh Token};
      \node[normal] (newtoken) at (3.95,-2.55) {返回新 Access Token\\更新 Cookie};

      \draw[arrow] (login) -- (verify);
      \draw[arrow] (verify) -- (issue);
      \draw[arrow] (issue) -- (save);
      \draw[arrow] (save) -- (cookie);
      \draw[arrow] (request) -- (auth);
      \draw[arrow] (auth) -- (valid);
      \draw[arrow] (valid) -- (biz);
      \draw[arrow] (auth.south) -- (expired.north);
      \draw[arrow] (expired.west) -- (refresh.east);
      \draw[arrow] (refresh.south) -- (lock.north);
      \draw[arrow] (lock) -- (hash);
      \draw[arrow] (hash) -- (rotate);
      \draw[arrow] (rotate) -- (newtoken);
    \end{tikzpicture}
  }
  \caption{用户登录与令牌刷新安全流程}\label{fig:implementation-auth-token-security-flow}
\end{figure}

如\cref{fig:implementation-auth-token-security-flow} 所示，刷新流程不是简单续期，而是令牌轮换。BIZ 使用 Redis 对同一个 Refresh Token 设置短时刷新锁，防止多个并发请求同时刷新；刷新成功后签发新的 Access Token 和 Refresh Token，并吊销旧 Refresh Token。Refresh Token 通过 HttpOnly Cookie 保存，前端 JavaScript 无法直接读取，降低脚本窃取风险。用户登出时，系统会将未过期 Access Token 的 jti 加入 Redis 黑名单，并吊销对应 Refresh Token。

项目权限由后端统一校验。所有项目相关接口都必须携带用户身份和项目 ID，Core BIZ 检查项目归属后才允许读取、修改或推进阶段。前端只负责隐藏不可用按钮，不能作为权限判断来源。WebSocket 连接也不直接使用访问令牌拼接 URL，而是先通过鉴权 HTTP 接口创建短时 Ticket，再用 Ticket 建立连接，降低长连接 URL 暴露后的风险。

对象存储不向前端暴露长期凭证。前端上传和下载图像使用预签名 URL，Reasoning 上传产物使用短时上传策略，Summary 下载报告源文件也使用预签名 URL。这样，即使前端或 Activity 服务拿到某个 URL，其权限也被限制在具体对象、具体前缀和有限时间内。

密码安全采用服务端哈希存储。注册和重置密码时，BIZ 会进行密码长度和复杂度校验，并使用 Bcrypt 生成密码哈希后写入数据库；登录时通过 Bcrypt 比对哈希，不保存明文密码。邮箱验证码按登录、注册和重置密码等场景区分，并通过 Redis 与签名密钥完成验证和消费，避免验证码跨场景复用。

WebSocket 安全采用“鉴权接口创建票据 + Scope 隔离”的方式。前端不能直接凭项目 ID 建立长连接，而是先调用受鉴权保护的 \texttt{/socket/ticket} 接口，由 BIZ 校验用户和项目权限后生成短时 Ticket。建立连接时需要同时携带项目 ID、Scope 和 ticket，API 再调用 BIZ 校验票据。Scope 将图像资产、检测任务和报告事件分开，避免一个页面接收不属于当前阶段的事件。

Agent 输出也需要安全性和可靠性约束。分类输出必须被解析为受控任务代码，非法任务代码会被拦截；总结输出必须被压缩为合法 JSON 对象；项目报告输出还需要校验图像组数量和图像数量是否与源数据一致。Reasoning 服务启动时校验模型文件，避免部署中只存在 Git LFS 指针而非真实模型权重。通过这些约束，系统将外部 Agent 和模型的非确定性限制在可处理范围内。

\section{稳定性设计与实现}
\label{sec:implementation-stability-design-and-implementation}

系统稳定性主要体现在日志、异常处理、容量控制、超时控制、幂等重试、产物校验和清理机制上。日志方面，各服务按 HTTP、RPC、MQ、WS、CRON、CLS、RSN、SMR 等范围输出结构化日志，并保留请求 ID、任务 UUID、图像 UUID、状态、耗时和错误信息。Docker Compose 中为五个后端服务配置日志采集 sidecar，将容器日志按 PSM 写入 \texttt{/root/logs/\{psm\}/log}，并按文件大小和最大文件数量进行轮转。

容量控制方面，Classification、Reasoning、Summary 和 Detection Worker 都有最大并发配置。Agent 调用和模型推理属于高耗时任务，如果不限制并发，容易压垮外部 Agent 平台、Python 模型进程、对象存储和数据库。通过并发上限和任务队列，系统在吞吐量与稳定性之间取得平衡。

超时控制方面，Agent 请求、Reasoning 总任务、产物下载、产物上传、图像上传待确认、Socket ticket 和预签名 URL 都有独立 TTL 或 timeout。Reasoning 服务在真正执行时再向 BIZ 获取原图 URL 和产物上传授权，避免长队列导致授权过期。上传中图像超时 CronJob 会将长期未上报的图像标记为失败并推送前端，避免页面长期停留在上传中。

数据一致性方面，系统将状态推进集中在 BIZ 服务中执行，并使用数据库事务包裹关键更新。Activity 服务完成任务后只通过回调上报结果，不能直接写库；BIZ 收到回调后校验任务是否仍然可访问、状态是否可替换、产物是否完整，再更新状态并发布事件。对于 Reasoning 产物，系统要求全部必要产物上传成功后才允许子任务成功，避免前端结果页出现缺图或对象不存在的问题。

异常恢复方面，检测和报告均支持失败状态和重试。检测失败后，前端可以重新启动失败任务；报告失败后，前端显示失败重试按钮，由 BIZ 重新创建报告任务并调用 Summary 服务。MinIO 脏对象清理 CronJob 负责删除失败图像、失败子任务、失败报告和无数据库引用对象对应的文件夹，降低对象存储长期膨胀风险。清理前设置保护窗口，避免并发执行中的对象被误删。

可观测性方面，工作台提供项目数量、图像数量、图像组数量、中间产物数量、检测执行次数、报告生成次数和对象存储容量概览。虽然它不是完整监控系统，但能让用户和维护者从账号维度观察系统使用情况和存储占用。结合服务日志、RocketMQ 消费日志、WebSocket 广播日志和 CronJob 日志，可以定位大部分业务链路问题。

\section{部署设计与实现}
\label{sec:implementation-deployment-implementation}

系统采用 Docker Compose 部署。Compose 文件中定义 RocketMQ NameServer、Broker、Topic 初始化服务、三个 Activity 服务、Core BIZ、Core API、Core Web 以及日志采集 sidecar。所有服务挂载到统一 Docker 网络中，服务之间通过容器服务名访问。Core API 对外暴露 HTTP 端口，Core Web 使用 Nginx 托管前端静态资源并映射到 80 端口，Reasoning 服务通过只读挂载方式加载服务器上的模型目录。

部署顺序由 \texttt{depends\_on} 和健康检查控制。RocketMQ Broker 健康后执行 Topic 初始化，Activity 服务、BIZ 和 API 均依赖 RocketMQ 初始化完成。BIZ 还依赖 Classification、Reasoning 和 Summary 服务启动，以保证检测 Worker 调用目标可用。这样的启动顺序降低了服务刚启动时因依赖未就绪导致的异常。

环境变量集中描述服务地址、MySQL、Redis、MinIO、RocketMQ、JWT、邮箱验证码、Socket ticket、对象存储 TTL、Agent Token、Agent Bot ID、任务超时和并发配置。与把配置写死在代码中相比，环境变量便于本地和服务器切换，也便于毕业设计系统在不同机器上复现部署。

Reasoning 模型文件没有打入 Docker 镜像，而是通过 \texttt{/root/models:/app/models:ro} 方式挂载。这样可以避免模型权重导致镜像体积过大，也便于更新模型文件。服务启动时会检查模型文件是否存在且不是 Git LFS 指针，提前暴露部署问题。

日志部署采用两层机制。容器自身使用 Docker \texttt{json-file} 驱动进行基础轮转，日志采集 sidecar 通过 Docker socket 读取目标容器日志，并按 PSM 写入宿主机目录。该方案比在每个业务容器中接入复杂日志组件更轻量，适合当前项目规模，也便于服务器直接查看各服务日志。

\section{本章小结}
\label{sec:implementation-chapter-summary}

本章从工程角度说明了第四章技术方案的落地方式。系统通过 Core Web、Core API、Core BIZ、Activity Classification、Activity Reasoning、Activity Summary 和 ICW Common 七类模块，将前端交互、业务状态、Agent 调用、模型推理、报告生成和共享协议拆分到清晰边界中。BIZ 服务集中控制任务状态和数据库事务，Activity 服务只返回结构化结果，API 服务负责接入与实时推送，Web 服务负责阶段化展示和用户复核。

工程实现的主要工作量体现在长链路协同上：一次检测任务涉及 HTTP 请求、gRPC 调用、Worker 入队、分类 Agent、子任务创建、Python 模型执行、MinIO 产物上传、回调落库、RocketMQ 事件、WebSocket 推送和前端状态合并；一次项目报告任务还涉及项目级源数据组装、source.txt 存储、附件式 Agent 调用、报告 JSON 校验和报告展示。通过协议治理、状态机、对象存储、CronJob、日志、并发控制和 Docker 部署，系统将 Agent 能力转化为可运行、可追踪、可恢复的软件系统。
